\documentclass[12pt,a4paper]{amsart}

\usepackage{amsmath,amssymb,amsthm}
\usepackage{mathtools}
\usepackage[utf8]{inputenc}
\usepackage[T1]{fontenc}
\usepackage[ngerman]{babel}
\usepackage{hyperref}
\usepackage{enumitem,booktabs,array}
\usepackage{geometry}
\geometry{margin=2.5cm}

% -------------------------------------------------------
% Theorem-Umgebungen (Deutsch)
% -------------------------------------------------------
\newtheorem{theorem}{Satz}[section]
\newtheorem{lemma}[theorem]{Lemma}
\newtheorem{korollar}[theorem]{Korollar}
\newtheorem{proposition}[theorem]{Proposition}
\theoremstyle{definition}
\newtheorem{definition}[theorem]{Definition}
\newtheorem{beispiel}[theorem]{Beispiel}
\theoremstyle{remark}
\newtheorem{bemerkung}[theorem]{Bemerkung}
\newtheorem{vermutung}[theorem]{Vermutung}

% -------------------------------------------------------
% Abkürzungen
% -------------------------------------------------------
\newcommand{\N}{\mathbb{N}}
\newcommand{\Z}{\mathbb{Z}}
\newcommand{\Q}{\mathbb{Q}}
\newcommand{\R}{\mathbb{R}}
\newcommand{\C}{\mathbb{C}}
\newcommand{\Ric}{\mathrm{Ric}}
\newcommand{\Rm}{\mathrm{Rm}}
\newcommand{\vol}{\mathrm{vol}}
\newcommand{\diam}{\mathrm{diam}}
\newcommand{\inj}{\mathrm{inj}}
\DeclareMathOperator{\tr}{tr}
\DeclareMathOperator{\divop}{div}

% -------------------------------------------------------
\begin{document}
% -------------------------------------------------------

\title{Ricci-Fluss und die Poincaré-Vermutung:\\
       Hamiltons Programm, Perelmans Beweis\\
       und der Geometrisierungssatz}

\author{Michael Fuhrmann}
\address{Specialist Mathematics Project, Build 147}
\email{specialist-maths@localhost}
\date{\today}

\begin{abstract}
Der Ricci-Fluss, 1982 von Richard Hamilton eingeführt, ist eine geometrische
Evolutionsgleichung
\[
  \frac{\partial g}{\partial t} = -2\,\Ric(g),
\]
die eine Riemannsche Metrik in Richtung ihrer Ricci-Krümmung deformiert.
Hamilton nutzte ihn, um zu beweisen, dass jede kompakte Drei-Mannigfaltigkeit
mit positiver Ricci-Krümmung diffeomorph zu einem sphärischen Raumformular
ist, und begann ein Programm zur Beweisführung von Thurstons
Geometrisierungsvermutung.  In den Jahren 2002--2003 vollendete Grigori
Perelman dieses Programm, indem er die $\mathcal{F}$- und
$\mathcal{W}$-Entropiefunktionale, das reduzierte Volumen und die Technik
des Ricci-Flusses mit Chirurgie einführte.  Als Korollar erhielt er die
Poincaré-Vermutung: Jede einfach zusammenhängende, geschlossene, orientierbare
Drei-Mannigfaltigkeit ist homöomorph zu $S^3$.  Beide Resultate sind
vollständig bewiesene Sätze.  Wir geben einen Überblick über die
wesentlichen Definitionen, Schlüsselabschätzungen, das Chirurgieverfahren
und die logische Struktur von Perelmans Argument, und stellen die
dreidimensionale Situation dem vierdimensionalen Fall gegenüber, in dem die
glatte Poincaré-Vermutung offen bleibt.
\end{abstract}

\maketitle
\tableofcontents

% ================================================================
\section{Einleitung}
% ================================================================

Die Poincaré-Vermutung, 1904 von Henri Poincaré formuliert, fragt, ob jede
einfach zusammenhängende, geschlossene $3$-Mannigfaltigkeit homöomorph zur
Drei-Sphäre $S^3$ ist.  Fast ein Jahrhundert lang widerstand sie allen
Beweisversuchen und wurde zum berühmtesten offenen Problem der Topologie.

Richard Hamilton~\cite{Hamilton1982} schlug 1982 eine Strategie vor, die auf
\emph{geometrischer Analysis} basiert: Deformiere eine beliebige Riemannsche
Metrik durch den \emph{Ricci-Fluss}, bis sie gegen eine runde Metrik
konvergiert.  In den folgenden zwei Jahrzehnten entwickelte er tiefe
Abschätzungen, klassifizierte Singularitäten und bewies mehrere grundlegende
Sätze.  Das Problem der Singularitäten in Dimension drei und höher schien
jedoch einen vollständigen Beweis zu verhindern.

Im Oktober 2002 stellte Grigori Perelman den ersten von drei Preprints auf
dem arXiv ein~\cite{Perelman2002,Perelman2003a,Perelman2003b}, der Hamiltons
Programm vervollständigte.  Perelman führte völlig neue monotone Größen ein
($\mathcal{F}$-Entropie, $\mathcal{W}$-Entropie, reduziertes Volumen),
bewies den $\kappa$-Nicht-Kollaps-Satz und konstruierte einen kanonischen
Ricci-Fluss mit Chirurgie, der alle endlichzeitigen Singularitäten übersteht.
Daraus leitete er sowohl die \emph{Geometrisierungsvermutung} (Thurston 1982)
als auch, als Spezialfall, die Poincaré-Vermutung ab.

Perelman wurde 2006 für diese Arbeit mit der Fields-Medaille ausgezeichnet;
er lehnte sie ab.  Vollständige, begutachtete Darstellungen lieferten
anschließend Morgan--Tian~\cite{MorganTian2007} und Cao--Zhu~\cite{CaoZhu2006}.

\subsection*{Gliederung}
Abschnitt~\ref{sec:riem} rekapituliert die Riemannsche Geometrie.
Abschnitt~\ref{sec:rf} definiert den Ricci-Fluss und den normierten Ricci-Fluss.
Abschnitt~\ref{sec:hamilton} stellt Hamiltons wichtigste Ergebnisse vor.
Abschnitt~\ref{sec:chirurgie} führt den Ricci-Fluss mit Chirurgie ein.
Abschnitt~\ref{sec:perelman} beschreibt Perelmans Entropiefunktionale
und Monotonieergebnisse.
Abschnitt~\ref{sec:poincare} formuliert und beweist (durch Referenz)
die Poincaré-Vermutung.
Abschnitt~\ref{sec:geom} behandelt den Geometrisierungssatz.
Abschnitt~\ref{sec:4d} diskutiert den vierdimensionalen Fall.
Abschnitt~\ref{sec:anwendungen} erwähnt weitere Anwendungen.


% ================================================================
\section{Riemannsche Geometrie -- Grundlagen}
\label{sec:riem}
% ================================================================

Sei $(M,g)$ eine glatte Riemannsche Mannigfaltigkeit der Dimension $n$.

\begin{definition}[Riemann-Krümmungstensor]
Der \emph{Riemann-Krümmungstensor} $\Rm$ ist der $(1,3)$-Tensor
\[
  \Rm(X,Y)Z = \nabla_X\nabla_Y Z - \nabla_Y\nabla_X Z - \nabla_{[X,Y]}Z,
\]
wobei $\nabla$ die Levi-Civita-Ableitung bezeichnet.  In lokalen
Koordinaten schreibt man $R^l{}_{ijk}$ für seine Komponenten.
\end{definition}

\begin{definition}[Ricci-Tensor]
Der \emph{Ricci-Tensor} ist die Kontraktion
\[
  R_{ij} = R^k{}_{ikj} = g^{kl}R_{kilj}.
\]
Er ist ein symmetrischer $(0,2)$-Tensor und misst grob, wie stark das Volumen
einer kleinen geodätischen Kugel vom euklidischen Volumen abweicht.
\end{definition}

\begin{definition}[Skalare Krümmung]
Die \emph{skalare Krümmung} ist die volle Spur
\[
  R = g^{ij}R_{ij} = \tr_g \Ric.
\]
\end{definition}

\begin{definition}[Schnittkrümmung]
Für eine Zweiebenenschicht $\sigma = \mathrm{span}\{X,Y\} \subset T_pM$:
\[
  K(\sigma) = \frac{g(\Rm(X,Y)Y,X)}{g(X,X)g(Y,Y)-g(X,Y)^2}.
\]
\end{definition}

\begin{definition}[Einstein-Mannigfaltigkeit]
$(M,g)$ heißt \emph{Einstein}, falls $\Ric = \lambda g$ für eine Konstante
$\lambda \in \R$.  Auf einer kompakten Einstein-Drei-Mannigfaltigkeit gilt
$\lambda > 0$, $= 0$ oder $< 0$ entsprechend positiver, verschwindender
oder negativer skalarer Krümmung.
\end{definition}

\begin{bemerkung}
In Dimension drei ist der volle Riemann-Krümmungstensor vollständig durch den
Ricci-Tensor bestimmt:
\[
  R_{ijkl} = g_{ik}R_{jl} - g_{il}R_{jk} - g_{jk}R_{il} + g_{jl}R_{ik}
             - \tfrac{1}{2}R(g_{ik}g_{jl} - g_{il}g_{jk}).
\]
Diese algebraische Tatsache spielt in Hamiltons Analysis eine Schlüsselrolle.
\end{bemerkung}

\begin{definition}[Gradient, Laplace-Operator, Divergenz]
Für eine Funktion $f$ und einen $(0,2)$-Tensor $T$:
\begin{align*}
  (\nabla f)^i &= g^{ij}\partial_j f, \\
  \Delta f &= g^{ij}\nabla_i\nabla_j f, \\
  (\divop T)_j &= g^{ik}\nabla_i T_{kj}.
\end{align*}
\end{definition}


% ================================================================
\section{Der Ricci-Fluss}
\label{sec:rf}
% ================================================================

\begin{definition}[Ricci-Fluss]
Sei $M$ eine kompakte glatte Mannigfaltigkeit und $g_0$ eine Riemannsche
Metrik auf $M$.  Der \emph{Ricci-Fluss} ist die einparametrige Familie von
Metriken $g(t)$, die
\begin{equation}
  \label{eq:rf}
  \frac{\partial}{\partial t} g(t) = -2\,\Ric(g(t)), \qquad g(0) = g_0
\end{equation}
erfüllt.
\end{definition}

Der Faktor $2$ ist konventionell; er stellt sicher, dass der Fluss auf dem
flachen Raum mit der Wärmeleitungsgleichung übereinstimmt.

\begin{bemerkung}[Heuristik]
Der Ricci-Fluss ist eine Art nichtlineare Wärmeleitungsgleichung für die Metrik.
Wo die Krümmung positiv ist (die Mannigfaltigkeit ist „zu stark gekrümmt"),
zieht sich die Metrik zusammen; wo die Krümmung negativ ist, dehnt sie sich aus.
Auf diese Weise tendiert der Fluss dazu, die Krümmung zu homogenisieren.
\end{bemerkung}

\begin{theorem}[Kurzfristige Existenz und Eindeutigkeit; DeTurck 1983]
\label{thm:ste}
Für jede glatte kompakte Riemannsche Mannigfaltigkeit $(M,g_0)$ existiert
$T > 0$ und eine eindeutige glatte Lösung $g(t)$ von~\eqref{eq:rf} auf
$[0,T)$.
\end{theorem}

\begin{proof}[Beweisskizze]
Die Ricci-Fluss-Gleichung ist wegen der Diffeomorphismeninvarianz nicht
strikt parabolisch.  DeTurcks Trick besteht darin,~\eqref{eq:rf} mit einer
zeitabhängigen Familie von Diffeomorphismen $\varphi_t$ zu koppeln, sodass
der zurückgeholte Fluss $\tilde g(t) = \varphi_t^* g(t)$ eine strikt
parabolische Gleichung erfüllt.  Die Standardtheorie parabolischer PDGlen
liefert dann kurzfristige Existenz und Eindeutigkeit von $\tilde g(t)$,
woraus man $g(t)$ zurückgewinnt.
\end{proof}

\subsection{Evolutionsgleichungen}

Unter dem Ricci-Fluss entwickeln sich die geometrischen Größen wie folgt:
\begin{align}
  \frac{\partial}{\partial t} R_{ij} &= \Delta R_{ij} + 2 R_{kijl}R^{kl}
     - 2 R_{ik}R^k{}_j, \label{eq:ric-ev} \\
  \frac{\partial}{\partial t} R &= \Delta R + 2|\Ric|^2. \label{eq:scal-ev}
\end{align}
Gleichung~\eqref{eq:scal-ev} impliziert via Maximumprinzip, dass
$R(x,t) \geq R_{\min}(0)$ für alle $t \geq 0$, d.\,h.\ das Infimum von $R$
ist monoton wachsend.

\subsection{Normierter Ricci-Fluss}

Um Konvergenz bei Mannigfaltigkeiten mit nicht-verschwindender
Euler-Charakteristik zu untersuchen, ist es zweckmäßig, die Metrik so
umzuskalieren, dass das Volumen erhalten bleibt.  Sei
$r(t) = \frac{\int_M R\,d\mu}{\vol(M,g(t))}$ die mittlere skalare Krümmung.
Der \emph{normierte Ricci-Fluss} lautet:
\begin{equation}
  \label{eq:nrf}
  \frac{\partial}{\partial t} g = -2\Ric + \frac{2}{n} r\, g.
\end{equation}
In Dimension $3$ ist der normierte Fluss äquivalent zum unnormierten Fluss
bis auf eine Zeitumparametrisierung und eine homothetische Umskalierung.

\subsection{Ricci-Solitone}

\begin{definition}[Ricci-Soliton]
Eine Riemannsche Metrik $g$ heißt \emph{Ricci-Soliton}, falls
\[
  \Ric + \nabla^2 f = \lambda g
\]
für eine glatte Funktion $f$ (das \emph{Potential}) und eine Konstante
$\lambda \in \R$.  Das Soliton heißt \emph{schrumpfend} ($\lambda > 0$),
\emph{stationär} ($\lambda = 0$) oder \emph{expandierend} ($\lambda < 0$).
\end{definition}

Ein Ricci-Soliton ist ein Fixpunkt des normierten Ricci-Flusses modulo
Diffeomorphismen und Skalierung.  Die runde Sphäre $S^n$ mit der
Standardmetrik ist ein schrumpfendes Ricci-Soliton mit $f = \mathrm{const}$
(eine Ricci--Einstein-Metrik).  Das Gauß-Soliton auf $\R^n$ ist ein
stationäres Soliton mit $f(x) = |x|^2/4$.


% ================================================================
\section{Hamiltons Ergebnisse}
\label{sec:hamilton}
% ================================================================

\begin{theorem}[Hamilton 1982]
\label{thm:hamilton82}
Sei $(M^3, g_0)$ eine kompakte Drei-Mannigfaltigkeit mit positiver
Ricci-Krümmung, $\Ric(g_0) > 0$.  Dann existiert der normierte Ricci-Fluss
für alle Zeiten, und die Metrik $g(t)$ konvergiert im $C^\infty$-Sinn gegen
eine Metrik konstanter positiver Schnittkrümmung für $t\to\infty$.
Insbesondere ist $M$ diffeomorph zu einem sphärischen Raumformular
$S^3/\Gamma$.
\end{theorem}

\begin{proof}[Strategie]
Hamilton analysierte die Entwicklung des Krümmungstensors mithilfe des
Maximumprinzips.  Die Schlüsselschritte sind:
\begin{enumerate}[label=(\roman*)]
  \item Zeige, dass $\Ric > 0$ unter dem Fluss erhalten bleibt
    (Pinching-Lemma).
  \item Beweise die \emph{Pinching-Abschätzung}: Falls $\Ric \geq \delta R\,g$
    zur Zeit $t = 0$, verbessert sich diese untere Schranke im Laufe der Zeit.
  \item Nutze die Harnack-Ungleichung für die Krümmung, um räumliche
    Variationen zu kontrollieren.
  \item Leite exponentielle Konvergenz gegen eine Metrik konstanter
    Krümmung via normierten Fluss ab.
\end{enumerate}
Die Annahme positiver Ricci-Krümmung ist wesentlich; sie verhindert die
Bildung langer dünner „Hälse" ($S^2 \times \R$), die sonst die Konvergenz
blockieren würden.
\end{proof}

\begin{korollar}
Jede kompakte Drei-Mannigfaltigkeit mit positiver Ricci-Krümmung hat eine
endliche Fundamentalgruppe.
\end{korollar}

\subsection{Hamiltons Entropie}

Hamilton führte das \emph{Entropiefunktional}
\[
  \mathcal{E}(g) = \int_M R \log R\, d\mu
\]
ein und zeigte, dass es auf Flächen unter dem normierten Ricci-Fluss monoton
wächst -- eine Rolle, die Perelmans $\mathcal{W}$-Entropie
(Abschnitt~\ref{sec:perelman}) verallgemeinert.

\subsection{Singularitätenanalyse}

In Dimension $\geq 3$ kann der Ricci-Fluss endlichzeitige Singularitäten
entwickeln.  Hamilton klassifizierte sie in:
\begin{itemize}
  \item \textbf{Typ I}: $|\Rm|(x,t) \leq \frac{C}{T-t}$ für $t \nearrow T$.
    Dazu gehört das runde Kollabieren einer Sphäre.
  \item \textbf{Typ II}: $|\Rm|$ wächst schneller als $(T-t)^{-1}$.
    Dazu gehört das Zigarren-Soliton.
  \item \textbf{Typ III}: Singularitäten für $t \to +\infty$ auf unsterblichen
    Lösungen.
\end{itemize}
Für Drei-Mannigfaltigkeiten mit positiver Ricci-Krümmung treten nur
Typ-I-Singularitäten auf (rundes Kollabieren), die durch
Satz~\ref{thm:hamilton82} aufgelöst werden.  Im allgemeinen Fall ohne
Positivität der Ricci-Krümmung können Hals-Einschnürungs-Singularitäten
(Typ I) entstehen; diese erfordern Chirurgie.


% ================================================================
\section{Ricci-Fluss mit Chirurgie}
\label{sec:chirurgie}
% ================================================================

\subsection{Hals-Einschnürungs-Singularitäten}

Die wichtigste endlichzeitige Singularität in Dimension drei ist die
\emph{Halseinschnürung} (neck-pinch): Eine Region diffeomorph zu
$S^2 \times (-1,1)$ entwickelt unbeschränkte Krümmung, wenn der
$S^2$-Faktor kollabiert.

\begin{definition}[$\varepsilon$-Hals]
Ein \emph{$\varepsilon$-Hals} in $(M,g)$ um $x$ ist eine Region
$U \subset M$ mit einem Diffeomorphismus $\psi: S^2 \times (-\varepsilon^{-1},
\varepsilon^{-1}) \to U$, sodass
\[
  \bigl\|\, r^{-2}\,\psi^*g - g_{S^2} \oplus ds^2 \,\bigr\|_{C^{[\varepsilon^{-1}]}} < \varepsilon,
\]
wobei $r$ der Radius des $S^2$-Faktors ist.
\end{definition}

\subsection{Kanonischer-Nachbarschafts-Satz}

Perelmans \emph{Satz über kanonische Nachbarschaften} besagt: Jedes $(x,t)$
mit hinreichend großer Krümmung $|\Rm|(x,t) \geq r^{-2}$ hat eine
Nachbarschaft (nach Umskalierung mit $|\Rm|(x,t)$), die $\varepsilon$-nah an
einem Standardmodell liegt:
\begin{itemize}
  \item ein starker $\varepsilon$-Hals,
  \item eine $\varepsilon$-Kappe (Topologie eines Halbraums oder
    Projektiv\-ebenen-Caps), oder
  \item eine kompakte $\varepsilon$-runde Mannigfaltigkeit.
\end{itemize}

\subsection{Chirurgieverfahren}

Zur Singularitätszeit $T$ führt man Chirurgie an allen Hälsen durch:
\begin{enumerate}
  \item Schneide $M$ entlang des Halses $S^2$ auf.
  \item Verschließe jede entstehende Randkomponente $S^2$ mit einer
    Standard-Halbkugel (\emph{Standardkappe}).
  \item Starte den Fluss mit der modifizierten Metrik $g^+(T)$ neu.
\end{enumerate}

\begin{theorem}[Perelman 2003; Existenz des chirurgischen Flusses]
\label{thm:chirurgie}
Es gibt Konstanten $r, \delta, \kappa > 0$, die nur von der Anfangsmetrik
abhängen, sodass der Ricci-Fluss mit $(r,\delta)$-Abschnitt für alle Zeiten
$t \in [0,\infty)$ existiert, mit nur endlich vielen Chirurgien auf jedem
endlichen Zeitintervall.
\end{theorem}

Ein entscheidender Punkt ist, dass die Chirurgien auf einer sorgfältig
gewählten Skala $h$ durchgeführt werden, die mit der Krümmungsskala $r$
vergleichbar ist, sodass die geometrische Qualität der Metrik nach der
Chirurgie kontrolliert ist.


% ================================================================
\section{Perelmans Beweis: Entropie, Nicht-Kollaps und reduziertes Volumen}
\label{sec:perelman}
% ================================================================

\subsection{Das $\mathcal{F}$-Entropiefunktional}

\begin{definition}[Perelmans $\mathcal{F}$-Funktional]
Sei $f: M \to \R$ eine glatte Funktion mit $\int_M e^{-f}\,d\mu = 1$.
Definiere
\begin{equation}
  \mathcal{F}(g,f) = \int_M \bigl(R + |\nabla f|^2\bigr) e^{-f}\,d\mu.
\end{equation}
\end{definition}

\begin{theorem}[Monotonie von $\mathcal{F}$; Perelman 2002]
\label{thm:F-mono}
Unter dem \emph{modifizierten Ricci-Fluss}
\[
  \partial_t g_{ij} = -2 R_{ij}, \qquad
  \partial_t f = -\Delta f - R + |\nabla f|^2,
\]
ist das $\mathcal{F}$-Funktional monoton wachsend:
\begin{equation}
  \frac{d}{dt}\mathcal{F}(g(t),f(t))
  = 2\int_M \bigl|R_{ij} + \nabla_i\nabla_j f\bigr|^2 e^{-f}\,d\mu \geq 0.
\end{equation}
Gleichheit gilt genau dann, wenn $(g,f)$ ein Gradient-stationäres
Ricci-Soliton ist.
\end{theorem}

\begin{proof}[Skizze]
Direkte Rechnung mithilfe der Evolutionsgleichungen für $g$, $d\mu$, $R$
und $|\nabla f|^2$.  Partielle Integration eliminiert Randterme auf einer
kompakten Mannigfaltigkeit.
\end{proof}

\subsection{Das $\mathcal{W}$-Entropiefunktional}

\begin{definition}[Perelmans $\mathcal{W}$-Funktional]
Für $\tau > 0$ und $(g, f)$ mit $\int_M (4\pi\tau)^{-n/2}e^{-f}\,d\mu = 1$:
\begin{equation}
  \mathcal{W}(g,f,\tau)
  = \int_M \bigl[\tau(R + |\nabla f|^2) + f - n\bigr]
    (4\pi\tau)^{-n/2}e^{-f}\,d\mu.
\end{equation}
\end{definition}

\begin{definition}[Perelmans $\mu$-Entropie]
\[
  \mu(g,\tau) = \inf \bigl\{\mathcal{W}(g,f,\tau) :
    \int_M (4\pi\tau)^{-n/2}e^{-f}\,d\mu = 1 \bigr\}.
\]
\end{definition}

\begin{theorem}[Monotonie von $\mathcal{W}$; Perelman 2002]
\label{thm:W-mono}
Unter dem Ricci-Fluss gekoppelt mit
\[
  \partial_t f = -\Delta f - R + \frac{n}{2\tau}, \qquad
  \frac{d\tau}{dt} = -1,
\]
ist das $\mathcal{W}$-Funktional monoton wachsend.  Insbesondere ist
$t \mapsto \mu(g(t), T-t)$ für $t < T$ monoton wachsend.
\end{theorem}

\subsection{$\kappa$-Nicht-Kollaps}

\begin{theorem}[$\kappa$-Nicht-Kollaps-Satz; Perelman 2002]
\label{thm:kappa}
Sei $(M,g(t))$ ein Ricci-Fluss auf $[0,T]$.  Für jedes $\rho > 0$ existiert
$\kappa > 0$ (abhängig von $\rho$, $T$ und $g(0)$), sodass gilt: Falls
$|\Rm|(x,t) \leq \rho^{-2}$ auf der Kugel $B(x,t,\rho)$, dann
\[
  \vol_t B(x,t,\rho) \geq \kappa \rho^n.
\]
\end{theorem}

\begin{proof}[Strategie]
Die untere Schranke für das $\mathcal{W}$-Entropiefunktional liefert eine
untere Schranke für $\mu(g(t),\rho^2)$.  Eine kollabierende Kugel würde
$\mathcal{W}$ sehr negativ machen, was dem Monotonieresultat~\ref{thm:W-mono}
widerspricht.
\end{proof}

Der $\kappa$-Nicht-Kollaps-Satz ist das zentrale Hilfsmittel: Er verhindert
die Entstehung zigarrenartiger Singularitäten, die zuvor Hamiltons Programm
blockierten, und stellt sicher, dass Regionen hoher Krümmung die
kanonische Nachbarschaftsstruktur aufweisen, die für die Chirurgie
erforderlich ist.

\subsection{Reduziertes Volumen und reduzierte Länge}

Perelman definierte die \emph{$\mathcal{L}$-Geodäte} und die
\emph{reduzierte Länge}:
\begin{equation}
  l(q,\tau) = \frac{1}{2\sqrt{\tau}}
    \inf_\gamma \int_0^\tau \sqrt{s}\bigl(R(\gamma(s)) + |\dot\gamma(s)|^2\bigr)ds,
\end{equation}
wobei das Infimum über alle Kurven $\gamma$ von $(p,0)$ nach $(q,\tau)$
(rückwärts in der Zeit durchlaufen) genommen wird.  Das
\emph{reduzierte Volumen} ist
\begin{equation}
  \tilde V(\tau) = \int_M (4\pi\tau)^{-n/2} e^{-l(q,\tau)}\,d\vol_{g(\tau)}.
\end{equation}

\begin{theorem}[Monotonie des reduzierten Volumens; Perelman 2002]
Unter dem Ricci-Fluss ist $\tilde V(\tau)$ monoton fallend in $\tau$,
und $\tilde V(\tau) \leq 1$.  Falls $\tilde V(\tau_0) = 1$ für ein
$\tau_0 > 0$, dann ist $(M,g)$ isometrisch zum euklidischen Raum.
\end{theorem}


% ================================================================
\section{Die Poincaré-Vermutung}
\label{sec:poincare}
% ================================================================

\begin{theorem}[Poincaré-Vermutung; Perelman 2003]
\label{thm:poincare}
Jede einfach zusammenhängende, geschlossene (kompakt, ohne Rand),
orientierbare Drei-Mannigfaltigkeit ist homöomorph zur Drei-Sphäre $S^3$.
\end{theorem}

\begin{proof}[Beweis via Ricci-Fluss mit Chirurgie]
Sei $M$ eine einfach zusammenhängende, geschlossene Drei-Mannigfaltigkeit.
Versehe $M$ mit einer beliebigen glatten Metrik $g_0$ und starte den
Ricci-Fluss mit Chirurgie (Satz~\ref{thm:chirurgie}).

\medskip
\textbf{Schritt 1: Der Fluss erlischt in endlicher Zeit.}
Da $\pi_1(M) = 0$, ist insbesondere $\pi_2(M) = 0$ (nach Hurewicz), und die
$3$-Mannigfaltigkeit ist prim.  Perelman~\cite{Perelman2003b} (und unabhängig
Colding--Minicozzi) zeigte, dass für einfach zusammenhängende geschlossene
$3$-Mannigfaltigkeiten der Ricci-Fluss mit Chirurgie in endlicher Zeit
$T < \infty$ erlischt: Die Mannigfaltigkeit wird nach endlich vielen
Chirurgien vollständig zerlegt, wobei jedes Stück rund kollabiert.

\medskip
\textbf{Schritt 2: Keine nicht-triviale Topologie überlebt.}
Kurz vor dem Erlöschen ist die Metrik $\varepsilon$-nah an einer runden
Metrik auf $S^3$ oder einem sphärischen Raumformular $S^3/\Gamma$.
Da $\pi_1(M) = 0$, kann kein nicht-triviales $\Gamma$ auftreten, sodass
nur $S^3$ möglich ist.

\medskip
\textbf{Schritt 3: Schlussfolgerung.}
Alle Chirurgien ersetzen $M$ durch eine topologisch äquivalente Mannigfaltigkeit
(nach der kanonischen Nachbarschaftsstruktur werden Chirurgien entlang
$S^2$-Hälsen durchgeführt und Standardkappen angesetzt, die topologische
Bälle sind).  Daher ist die ursprüngliche Mannigfaltigkeit $M$ homöomorph
zu $S^3$.
\end{proof}

\begin{bemerkung}
Die vollständige Verifikation von Perelmans Argument wurde von
Morgan--Tian~\cite{MorganTian2007} und Cao--Zhu~\cite{CaoZhu2006} unabhängig
voneinander durchgeführt.  Beide bestätigten alle Details des Beweises.
In der mathematischen Gemeinschaft besteht nun vollständiger Konsens darüber,
dass Satz~\ref{thm:poincare} bewiesen ist.
\end{bemerkung}

\begin{bemerkung}[Historischer Kontext]
Die Poincaré-Vermutung war eines der sieben Millennium-Probleme des Clay
Mathematics Institute (CMI), die jeweils mit einem Preis von
\$1{.}000{.}000 ausgezeichnet sind.  Perelman erhielt den Preis 2010, lehnte
ihn jedoch ab -- ebenso wie die Fields-Medaille 2006.
\end{bemerkung}


% ================================================================
\section{Die Geometrisierungsvermutung}
\label{sec:geom}
% ================================================================

\begin{definition}[Thurstons acht Geometrien]
\label{def:geom}
Eine dreidimensionale Geometrie $(X,G)$ besteht aus einer einfach
zusammenhängenden Drei-Mannigfaltigkeit $X$ und einer transitiven
Isometriegruppe $G$, deren Stabilisator eines Punktes kompakt ist.
Es gibt genau acht solche maximale Geometrien:
\[
  S^3,\quad \R^3,\quad H^3,\quad S^2\times\R,\quad H^2\times\R,\quad
  \widetilde{\mathrm{SL}_2(\R)},\quad \mathrm{Nil},\quad \mathrm{Sol}.
\]
\end{definition}

\begin{definition}[Geometrische Zerlegung]
Eine geschlossene, orientierbare Drei-Mannigfaltigkeit $M$ heißt
\emph{geometrisch}, wenn sie auf einer der acht Geometrien modelliert ist.
Die \emph{JSJ-Zerlegung} zerschneidet $M$ entlang inkompressibler Tori in
atoroidale Stücke; zusammen mit der Prim-Zerlegung entlang $S^2$ erhält man
die \emph{Geometrisierung}.
\end{definition}

\begin{theorem}[Geometrisierungssatz; Thurston 1982, Perelman 2003]
\label{thm:geom}
Jede geschlossene, orientierbare Drei-Mannigfaltigkeit lässt sich entlang
endlich vieler disjunkter, inkompressibler Tori und Sphären in Stücke
zerlegen, von denen jedes eine vollständige, lokal homogene Riemannsche
Metrik besitzt, die auf einer der acht Thurstonischen Geometrien modelliert
ist.
\end{theorem}

\begin{proof}[Beweis via Ricci-Fluss mit Chirurgie]
Perelmans Ricci-Fluss mit Chirurgie liefert eine kanonische Zerlegung von
$M$ in prim Stücke.  Für die atoroidalen (hyperbolischen) Stücke bewies
Perelman die \emph{Dick-Dünn-Zerlegung}: Nach Umskalierung konvergiert der
dicke Teil gegen eine hyperbolische Mannigfaltigkeit (mittels Hamiltons
Kompaktheitssatz und den $\kappa$-Nicht-Kollaps-Abschätzungen), während der
dünne Teil eine Graph-Mannigfaltigkeit ist (und daher eine der
Seifert-gefaserten Geometrien trägt).  Der Graph-Mannigfaltigkeiten-Fall
wurde separat mithilfe von Ergebnissen von Shioya--Yamaguchi behandelt.
Zusammen liefern diese den vollständigen Geometrisierungssatz.
\end{proof}

\begin{korollar}
Die Poincaré-Vermutung (Satz~\ref{thm:poincare}) folgt aus dem
Geometrisierungssatz: Die einzige Geometrie, die mit einer einfach
zusammenhängenden geschlossenen Drei-Mannigfaltigkeit kompatibel ist, ist die
sphärische Geometrie $S^3$, also $M \cong S^3$.
\end{korollar}

\begin{bemerkung}
Der Geometrisierungssatz wurde erstmals 1982 von William
Thurston~\cite{Thurston1982} vermutet, basierend auf seinen tiefen Arbeiten
über hyperbolische Drei-Mannigfaltigkeiten (wofür er 1982 die Fields-Medaille
erhielt).  Thurston bewies die Vermutung für Haken-Mannigfaltigkeiten; der
allgemeine Fall erforderte Perelmans Ricci-Fluss-Methoden.
\end{bemerkung}


% ================================================================
\section{Der vierdimensionale Fall}
\label{sec:4d}
% ================================================================

Die Situation in Dimension vier unterscheidet sich grundlegend von Dimension
drei.

\subsection{Topologische Poincaré-Vermutung in Dimension 4}

\begin{theorem}[Freedman 1982]
\label{thm:freedman}
Jede einfach zusammenhängende, geschlossene topologische Vier-Mannigfaltigkeit
mit dem Homotopietyp von $S^4$ ist homöomorph zu $S^4$.
\end{theorem}

Freedmans Beweis~\cite{Freedman1982} verwendet die Technik der
\emph{Casson-Henkel} und ist spezifisch für Dimension vier.  Er beruht auf
dem Whitney-Trick, der topologisch in Dimension vier, aber nicht glatt
funktioniert.  Freedman erhielt 1986 die Fields-Medaille für diese Arbeit.

\subsection{Glatte Poincaré-Vermutung in Dimension 4 -- offenes Problem}

\begin{vermutung}[Glatte vierdimensionale Poincaré-Vermutung]
\label{verm:s4d}
Jede glatte Mannigfaltigkeit, die homöomorph zu $S^4$ ist, ist auch
diffeomorph zu $S^4$.
\end{vermutung}

Dies bleibt eines der wichtigsten offenen Probleme der geometrischen
Topologie.  Im Gegensatz zur topologischen Aussage
(Satz~\ref{thm:freedman}) ist das glatte Problem vollständig offen.

\begin{bemerkung}
In allen anderen Dimensionen $n \neq 4$ ist die glatte Poincaré-Vermutung
entschieden:
\begin{itemize}
  \item $n \leq 2$: klassisch (trivial für $n=1$).
  \item $n = 3$: Perelmans Satz (Satz~\ref{thm:poincare}).
  \item $n \geq 5$: Smale (1961) bewies, dass jede einfach zusammenhängende,
    geschlossene $n$-Mannigfaltigkeit, die homotopieäquivalent zu $S^n$ ist,
    für $n \geq 5$ diffeomorph zu $S^n$ ist
    (der $h$-Kobordismussatz).
\end{itemize}
Dimension vier ist außergewöhnlich, weil der Whitney-Trick topologisch, aber
nicht glatt funktioniert.  Donaldsonsche Eichtheorie und Seiberg--Witten-Theorie
haben Hindernisse geliefert: Auf $\R^4$ existieren exotische glatte Strukturen
(auf $\R^n$ für $n \neq 4$ nicht), und viele einfach zusammenhängende
$4$-Mannigfaltigkeiten tragen unendlich viele verschiedene glatte Strukturen.
Für $S^4$ selbst wurde weder eine exotische glatte Struktur gefunden noch
ausgeschlossen.
\end{bemerkung}

\subsection{Ricci-Fluss in Dimension 4}

Der Ricci-Fluss in Dimension vier stößt auf neue Herausforderungen:
\begin{itemize}
  \item Der Weyl-Krümmungstensor verschwindet in Dimension vier nicht, sodass
    der Riemann-Tensor nicht allein durch $\Ric$ bestimmt wird.
  \item Neue Singularitätstypen (z.\,B.\ Orbifold-Singularitäten) treten auf.
  \item Hamiltons Singularitätenklassifikation ist weniger vollständig.
\end{itemize}
Fortschritte umfassen Arbeiten von Chen--Zhu und Brendle über
Vier-Mannigfaltigkeiten mit punktweise $\frac{1}{4}$-gepinnter
Schnittkrümmung.


% ================================================================
\section{Anwendungen und Verallgemeinerungen}
\label{sec:anwendungen}
% ================================================================

\subsection{Mittlerer Krümmungsfluss}

Der \emph{mittlere Krümmungsfluss} (Mean Curvature Flow, MCF) bewegt eine
Hyperfläche $\Sigma_t \subset \R^{n+1}$ in Richtung ihres mittleren
Krümmungsvektors:
\[
  \frac{d}{dt}\mathbf{x} = H\,\mathbf{n},
\]
wobei $H$ die mittlere Krümmung und $\mathbf{n}$ die innere Einheitsnormale ist.
Analog zum Ricci-Fluss entwickelt der MCF Singularitäten
(Halseinschnürungen), die durch Chirurgie aufgelöst werden können.
Andrews, Brendle, Huisken und Ilmanen entwickelten einen MCF mit Chirurgie
analog zu Perelmans Konstruktion.

\subsection{Kähler--Ricci-Fluss}

Auf einer Kähler-Mannigfaltigkeit $(M, J, g)$ repräsentiert die
Ricci-Form $\rho = -\frac{i}{2\pi}\partial\bar\partial\log\det(g_{i\bar j})$
die erste Chern-Klasse.  Der \emph{Kähler--Ricci-Fluss} lautet:
\[
  \frac{\partial \omega}{\partial t} = -\Ric(\omega) + \lambda\omega, \qquad \lambda \in \{-1, 0, 1\},
\]
und wird bei der Suche nach kanonischen Kähler-Metriken (Calabi--Yau,
Kähler--Einstein) eingesetzt.  Die Yau--Tian--Donaldson-Vermutung (bewiesen
von Chen, Donaldson und Sun 2015) besagt, dass K-Stabilität äquivalent zur
Existenz einer Kähler--Einstein-Metrik ist, wobei der Kähler--Ricci-Fluss
eine zentrale Rolle spielt.

\subsection{Harmonischer Kartenfluss und Yang--Mills-Fluss}

Die für den Ricci-Fluss entwickelten Techniken haben Analoga in:
\begin{itemize}
  \item \textbf{Harmonischer Kartenfluss}: $\partial_t u = \tau(u)$, wobei
    $\tau$ das Spannungsfeld ist.  Eells--Sampson (1964) bewiesen
    Langzeitexistenz, wenn das Ziel nicht-positive Krümmung hat.
  \item \textbf{Yang--Mills-Fluss}: Gradientenfluss des Yang--Mills-Energie\-funktionals.
    Donaldson nutzte ihn zum Beweis des Donaldson--Uhlenbeck--Yau-Satzes über
    stabile Bündel.
\end{itemize}

\subsection{Thurstons Programm und niedrigdimensionale Topologie}

Der Geometrisierungssatz löst die \emph{virtuelle Haken-Vermutung} (bewiesen
von Agol 2012) und die \emph{virtuelle Faserungsvermutung}, die beide nun
Sätze sind.  Die Klassifikation von Drei-Mannigfaltigkeiten ist damit
im Prinzip vollständig: Jede geschlossene orientierbare Drei-Mannigfaltigkeit
lässt sich anhand der acht Modellgeometrien verstehen.


% ================================================================
\section*{Danksagung}
% ================================================================

Dieses Paper ist Teil des \emph{Specialist Mathematics Project}, einer
systematischen Untersuchung fortgeschrittener Mathematik von Algebra und
Analysis bis zu den Millennium-Problemen.  Die Darstellung folgt den
Ausführungen in Morgan--Tian~\cite{MorganTian2007},
Cao--Zhu~\cite{CaoZhu2006} und Kleiner--Lott~\cite{KleinerLott2008}.


% ================================================================
\begin{thebibliography}{99}
% ================================================================

\bibitem{Hamilton1982}
R.~S.~Hamilton,
\textit{Three-manifolds with positive Ricci curvature},
J.\ Differential Geom.\ \textbf{17} (1982), Nr.~2, 255--306.

\bibitem{Perelman2002}
G.~Perelman,
\textit{The entropy formula for the Ricci flow and its geometric applications},
arXiv:math/0211159 (2002).

\bibitem{Perelman2003a}
G.~Perelman,
\textit{Ricci flow with surgery on three-manifolds},
arXiv:math/0303109 (2003).

\bibitem{Perelman2003b}
G.~Perelman,
\textit{Finite extinction time for the solutions to the Ricci flow on certain three-manifolds},
arXiv:math/0307245 (2003).

\bibitem{MorganTian2007}
J.~Morgan und G.~Tian,
\textit{Ricci Flow and the Poincar\'{e} Conjecture},
Clay Mathematics Monographs, Bd.~3, AMS, Providence, 2007.

\bibitem{CaoZhu2006}
H.-D.~Cao und X.-P.~Zhu,
\textit{A complete proof of the Poincar\'{e} and geometrization conjectures ---
application of the Hamilton-Perelman theory of the Ricci flow},
Asian J.\ Math.\ \textbf{10} (2006), Nr.~2, 165--492.

\bibitem{KleinerLott2008}
B.~Kleiner und J.~Lott,
\textit{Notes on Perelman's papers},
Geom.\ Topol.\ \textbf{12} (2008), 2587--2855.

\bibitem{Thurston1982}
W.~P.~Thurston,
\textit{Three-dimensional manifolds, Kleinian groups and hyperbolic geometry},
Bull.\ Amer.\ Math.\ Soc.\ \textbf{6} (1982), 357--381.

\bibitem{Freedman1982}
M.~H.~Freedman,
\textit{The topology of four-dimensional manifolds},
J.\ Differential Geom.\ \textbf{17} (1982), 357--453.

\bibitem{Chow2004}
B.~Chow und D.~Knopf,
\textit{The Ricci Flow: An Introduction},
Mathematical Surveys and Monographs, Bd.~110, AMS, Providence, 2004.

\end{thebibliography}

\end{document}
